\chapter{Phân tích bài toán và thiết kế hệ thống}
\label{chap:design}

Chương này mô tả việc chuyển từ lý thuyết sang một hệ thống mô phỏng chạy được: lựa chọn dữ liệu, kiến trúc phần mềm, mô hình dữ liệu và luồng tấn công.

\section{Lựa chọn tập dữ liệu}

Tập Netflix Prize gốc hiện không còn được phân phối công khai (chính vì các quan ngại riêng tư nói trên). Vì vậy chúng em dùng tập MovieLens 100k~\cite{movielens} làm proxy - một tập dữ liệu công khai do nhóm GroupLens (Đại học Minnesota) phát hành, có cùng bản chất (đánh giá phim kèm thời điểm) và cùng tính thưa. Thông số tập dữ liệu được tóm tắt ở Bảng~\ref{tab:dataset}.

\begin{table}[!htb]
\centering
\caption{Thông số tập dữ liệu MovieLens 100k}
\label{tab:dataset}
\renewcommand{\arraystretch}{1.3}
\begin{tabularx}{0.8\linewidth}{X r}
\toprule
\textbf{Thuộc tính} & \textbf{Giá trị} \\
\midrule
Số người dùng & 943 \\
Số phim & 1.682 \\
Số lượt đánh giá & 100.000 \\
Thang điểm & 1 - 5 (số nguyên) \\
Khoảng thời gian & 09/1997 - 04/1998 \\
Độ thưa (sparsity) & 93{,}70\% \\
\bottomrule
\end{tabularx}
\end{table}

Độ thưa $93{,}70\%$ nghĩa là trong ma trận $943 \times 1682$ (người $\times$ phim), chỉ khoảng $6{,}3\%$ ô có dữ liệu - phần còn lại để trống. Chính sự trống trải này tạo ra các tổ hợp phim gần như duy nhất cho mỗi người, đúng tiền đề mà thuật toán khai thác. Bảng~\ref{tab:mapping} đối chiếu vai trò giữa hai tập dữ liệu.

\begin{table}[!htb]
\centering
\caption{Đối chiếu tập Netflix gốc và proxy MovieLens 100k}
\label{tab:mapping}
\renewcommand{\arraystretch}{1.3}
\begin{tabularx}{\linewidth}{p{4.2cm} X X}
\toprule
\textbf{Khái niệm} & \textbf{Netflix Prize} & \textbf{MovieLens 100k} \\
\midrule
Bản ghi một cá nhân & Một thuê bao & Một người dùng \\
Thuộc tính & (phim, điểm, ngày) & (phim, điểm, ngày) \\
Quy mô & $\sim$500.000 người & 943 người \\
Tính thưa & Áp đảo & 93{,}70\% \\
\bottomrule
\end{tabularx}
\end{table}

Do MovieLens nhỏ và dày hơn Netflix, các con số tuyệt đối sẽ khác đôi chút (đám đông ứng viên ít hơn), nhưng \textit{xu hướng} - vốn là điều bài báo khẳng định - được tái hiện trung thực.

\section{Kiến trúc hệ thống}

Hệ thống mô phỏng theo mô hình client-server ba lớp, minh họa ở Hình~\ref{fig:arch}.

\begin{figure}[!htb]
\centering
\begin{tikzpicture}[font=\small, node distance=0pt]
\tikzset{
  box/.style={draw, rounded corners, minimum width=3.2cm, minimum height=1.15cm, align=center, thick},
}
\node[box, fill=blue!8] (fe) at (0,0) {\textbf{Frontend}\\ React + Vite + Tailwind\\ \footnotesize Bảng điều khiển kẻ tấn công};
\node[box, fill=green!8] (be) at (5.4,0) {\textbf{Backend}\\ FastAPI (Python)\\ \footnotesize Engine Scoreboard-RH};
\node[box, fill=orange!8] (db) at (10.8,0) {\textbf{CSDL}\\ SQLite + SQLAlchemy\\ \footnotesize MovieLens 100k};
\draw[->, thick] (fe.east) -- node[above]{\footnotesize REST/JSON} node[below]{\footnotesize /api} (be.west);
\draw[->, thick] (be.east) -- node[above]{\footnotesize SQL} node[below]{\footnotesize pandas} (db.west);
\end{tikzpicture}
\caption{Kiến trúc ba lớp của hệ thống mô phỏng tấn công giải ẩn danh}
\label{fig:arch}
\end{figure}

\renewcommand{\labelitemi}{$-$}
\begin{itemize}
	\item \textbf{Lớp dữ liệu (SQLite):} lưu ba bảng \texttt{users}, \texttt{movies}, \texttt{ratings} được nạp từ tệp gốc MovieLens. Truy cập qua SQLAlchemy.
	\item \textbf{Lớp backend (FastAPI):} chứa \textit{engine tấn công} - cài đặt thuật toán Scoreboard-RH bằng pandas/numpy. Toàn bộ 100.000 lượt đánh giá được nạp sẵn vào RAM khi khởi động và trọng số $wt(i)$ của mỗi phim được tính trước, nhờ đó mỗi lần tấn công chỉ mất vài chục mili-giây. Backend phơi bày REST API (xem Bảng~\ref{tab:api}).
	\item \textbf{Lớp frontend (React):} ``bảng điều khiển kẻ tấn công'' cho phép nhập tri thức bổ trợ và trực quan hóa quá trình lộ lọt thông tin.
\end{itemize}

\begin{table}[!htb]
\centering
\caption{Các điểm cuối (endpoint) REST của backend}
\label{tab:api}
\renewcommand{\arraystretch}{1.3}
\begin{tabularx}{\linewidth}{p{3.6cm} p{1.4cm} X}
\toprule
\textbf{Endpoint} & \textbf{Method} & \textbf{Chức năng} \\
\midrule
\texttt{/stats} & GET & Thống kê tập dữ liệu: độ thưa, phân phối điểm. \\
\texttt{/movies/\allowbreak search} & GET & Tìm phim theo tên, kèm support và trọng số $wt(i)$. \\
\texttt{/attack} & POST & Chạy thuật toán giải ẩn danh trên aux đã nhập. \\
\texttt{/auto\_aux} & GET & Sinh aux từ một nạn nhân có thật (có thể thêm nhiễu). \\
\bottomrule
\end{tabularx}
\end{table}

\section{Mô hình dữ liệu}

Ba bảng quan hệ được thiết kế tối giản, đủ cho bài toán:
\renewcommand{\labelitemi}{$-$}
\begin{itemize}
	\item Bảng \texttt{users} gồm các cột: id, age, gender, occupation, zip\_code - hồ sơ người dùng (đóng vai ``thông tin nhạy cảm'' có thể bị suy ra).
	\item Bảng \texttt{movies} gồm các cột: id, title, release\_date - danh mục phim.
	\item Bảng \texttt{ratings} gồm các cột: user\_id, movie\_id, rating, timestamp - bảng trung tâm, 100.000 dòng, là đối tượng chính của tấn công.
\end{itemize}

\section{Luồng tấn công}

Một phiên tấn công đi qua các bước sau (tương ứng với thao tác trên giao diện ở Chương~\ref{chap:exp}):
\begin{enumerate}[1.]
	\item Kẻ tấn công thu thập tri thức bổ trợ: một vài phim mà nạn nhân đã xem, kèm điểm và (tùy chọn) ngày xem gần đúng.
	\item Backend tính điểm cho mọi người dùng có chứa ít nhất một phim trong aux, theo công thức~(\ref{eq:score}).
	\item Backend tính eccentricity theo~(\ref{eq:ecc}) và so với ngưỡng $\phi$.
	\item Nếu vượt ngưỡng: kết luận danh tính, đồng thời truy xuất và phơi bày toàn bộ lịch sử đánh giá của nạn nhân - phần dữ liệu mà kẻ tấn công ban đầu không hề biết.
\end{enumerate}

Bước 4 chính là minh chứng trực quan cho ``rò rỉ riêng tư'': từ vài mẩu thông tin vụn vặt, kẻ tấn công thu được toàn bộ chân dung hành vi của nạn nhân.
